% Options for packages loaded elsewhere
% Options for packages loaded elsewhere
\PassOptionsToPackage{unicode}{hyperref}
\PassOptionsToPackage{hyphens}{url}
\PassOptionsToPackage{dvipsnames,svgnames,x11names}{xcolor}
%
\documentclass[
  a4paper,
  DIV=11,
  numbers=noendperiod]{scrartcl}
\usepackage{xcolor}
\usepackage[margin=35mm]{geometry}
\usepackage{amsmath,amssymb}
\setcounter{secnumdepth}{-\maxdimen} % remove section numbering
\usepackage{iftex}
\ifPDFTeX
  \usepackage[T1]{fontenc}
  \usepackage[utf8]{inputenc}
  \usepackage{textcomp} % provide euro and other symbols
\else % if luatex or xetex
  \usepackage{unicode-math} % this also loads fontspec
  \defaultfontfeatures{Scale=MatchLowercase}
  \defaultfontfeatures[\rmfamily]{Ligatures=TeX,Scale=1}
\fi
\usepackage{lmodern}
\ifPDFTeX\else
  % xetex/luatex font selection
\fi
% Use upquote if available, for straight quotes in verbatim environments
\IfFileExists{upquote.sty}{\usepackage{upquote}}{}
\IfFileExists{microtype.sty}{% use microtype if available
  \usepackage[]{microtype}
  \UseMicrotypeSet[protrusion]{basicmath} % disable protrusion for tt fonts
}{}
\makeatletter
\@ifundefined{KOMAClassName}{% if non-KOMA class
  \IfFileExists{parskip.sty}{%
    \usepackage{parskip}
  }{% else
    \setlength{\parindent}{0pt}
    \setlength{\parskip}{6pt plus 2pt minus 1pt}}
}{% if KOMA class
  \KOMAoptions{parskip=half}}
\makeatother
% Make \paragraph and \subparagraph free-standing
\makeatletter
\ifx\paragraph\undefined\else
  \let\oldparagraph\paragraph
  \renewcommand{\paragraph}{
    \@ifstar
      \xxxParagraphStar
      \xxxParagraphNoStar
  }
  \newcommand{\xxxParagraphStar}[1]{\oldparagraph*{#1}\mbox{}}
  \newcommand{\xxxParagraphNoStar}[1]{\oldparagraph{#1}\mbox{}}
\fi
\ifx\subparagraph\undefined\else
  \let\oldsubparagraph\subparagraph
  \renewcommand{\subparagraph}{
    \@ifstar
      \xxxSubParagraphStar
      \xxxSubParagraphNoStar
  }
  \newcommand{\xxxSubParagraphStar}[1]{\oldsubparagraph*{#1}\mbox{}}
  \newcommand{\xxxSubParagraphNoStar}[1]{\oldsubparagraph{#1}\mbox{}}
\fi
\makeatother


\usepackage{longtable,booktabs,array}
\usepackage{calc} % for calculating minipage widths
% Correct order of tables after \paragraph or \subparagraph
\usepackage{etoolbox}
\makeatletter
\patchcmd\longtable{\par}{\if@noskipsec\mbox{}\fi\par}{}{}
\makeatother
% Allow footnotes in longtable head/foot
\IfFileExists{footnotehyper.sty}{\usepackage{footnotehyper}}{\usepackage{footnote}}
\makesavenoteenv{longtable}
\usepackage{graphicx}
\makeatletter
\newsavebox\pandoc@box
\newcommand*\pandocbounded[1]{% scales image to fit in text height/width
  \sbox\pandoc@box{#1}%
  \Gscale@div\@tempa{\textheight}{\dimexpr\ht\pandoc@box+\dp\pandoc@box\relax}%
  \Gscale@div\@tempb{\linewidth}{\wd\pandoc@box}%
  \ifdim\@tempb\p@<\@tempa\p@\let\@tempa\@tempb\fi% select the smaller of both
  \ifdim\@tempa\p@<\p@\scalebox{\@tempa}{\usebox\pandoc@box}%
  \else\usebox{\pandoc@box}%
  \fi%
}
% Set default figure placement to htbp
\def\fps@figure{htbp}
\makeatother





\setlength{\emergencystretch}{3em} % prevent overfull lines

\providecommand{\tightlist}{%
  \setlength{\itemsep}{0pt}\setlength{\parskip}{0pt}}



 


\KOMAoption{captions}{tableheading}
\makeatletter
\@ifpackageloaded{caption}{}{\usepackage{caption}}
\AtBeginDocument{%
\ifdefined\contentsname
  \renewcommand*\contentsname{Table of contents}
\else
  \newcommand\contentsname{Table of contents}
\fi
\ifdefined\listfigurename
  \renewcommand*\listfigurename{List of Figures}
\else
  \newcommand\listfigurename{List of Figures}
\fi
\ifdefined\listtablename
  \renewcommand*\listtablename{List of Tables}
\else
  \newcommand\listtablename{List of Tables}
\fi
\ifdefined\figurename
  \renewcommand*\figurename{Figure}
\else
  \newcommand\figurename{Figure}
\fi
\ifdefined\tablename
  \renewcommand*\tablename{Table}
\else
  \newcommand\tablename{Table}
\fi
}
\@ifpackageloaded{float}{}{\usepackage{float}}
\floatstyle{ruled}
\@ifundefined{c@chapter}{\newfloat{codelisting}{h}{lop}}{\newfloat{codelisting}{h}{lop}[chapter]}
\floatname{codelisting}{Listing}
\newcommand*\listoflistings{\listof{codelisting}{List of Listings}}
\makeatother
\makeatletter
\makeatother
\makeatletter
\@ifpackageloaded{caption}{}{\usepackage{caption}}
\@ifpackageloaded{subcaption}{}{\usepackage{subcaption}}
\makeatother
\usepackage{bookmark}
\IfFileExists{xurl.sty}{\usepackage{xurl}}{} % add URL line breaks if available
\urlstyle{same}
\hypersetup{
  pdftitle={Skaler ve Vektör Kavramı},
  pdfauthor={Öğr. Gör. Oktay Cesur},
  colorlinks=true,
  linkcolor={blue},
  filecolor={Maroon},
  citecolor={Blue},
  urlcolor={Blue},
  pdfcreator={LaTeX via pandoc}}


\title{Skaler ve Vektör Kavramı}
\usepackage{etoolbox}
\makeatletter
\providecommand{\subtitle}[1]{% add subtitle to \maketitle
  \apptocmd{\@title}{\par {\large #1 \par}}{}{}
}
\makeatother
\subtitle{MATE 213 --- İlk Öğretim Bloğu}
\author{Öğr. Gör. Oktay Cesur}
\date{2026-07-23}
\begin{document}
\maketitle


\subsection{Bugünün Sorusu}\label{buguxfcnuxfcn-sorusu}

\begin{quote}
Birbiriyle ilişkili birden fazla sayıyı, sırasını ve anlamını
kaybetmeden nasıl tek bir nesne olarak ifade ederiz?
\end{quote}

\begin{itemize}
\tightlist
\item
  Tek bir sayı her durumda yeterli mi?\\
\item
  İlişkili sayılar beraber tutulabilir mi?
\item
  Beraber tuttuk diyelim, ölçekleyebilecek miyiz?
\end{itemize}

Bazen tek bir sayı bir durumu anlatmaya yetmez. Bir noktanın düzlemdeki
konumu iki sayı ister. Bir ürünün birkaç farklı özelliği olabilir. Bir
sistem de sıcaklık, basınç ve debi gibi birden fazla ölçümle
tanımlanabilir.

Bu tür durumlarda elimizde birbirinden kopuk sayılar değil, birlikte
anlam taşıyan bir sayı grubu vardır. Üstelik bu sayıların sırası da
önemlidir. Amacımız, bu bilgileri sırasını ve anlamını koruyarak tek bir
nesne halinde yazabilmek. Bunu önce koordinat vektörleriyle yapacağız.
Daha sonra bu nesneyi oluşturan sayılar ile onun tamamını büyütüp
küçültmek için kullandığımız sayıları birbirinden ayıracağız.

\begin{center}\rule{0.5\linewidth}{0.5pt}\end{center}

\subsection{İlk Örnek: Birlikte Anlamlı
Sayılar}\label{ilk-uxf6rnek-birlikte-anlamlux131-sayux131lar}

Bir ürünün üç özelliği:

\[
\text{fiyat: }40,\qquad
\text{ağırlık: }2,\qquad
\text{stok: }15
\]

\begin{itemize}
\tightlist
\item
  Değerler birbiriyle ilişkili,
\item
  sıra bozulursa anlam bozulur,
\item
  tek nesne olarak tutmak istiyoruz.
\end{itemize}

\[
v=\begin{bmatrix}40\\2\\15\end{bmatrix}
\]

Bu örnek, vektöre neden ihtiyaç duyduğumuzu geometrik ok çizmeden
gösterir. Başlangıçta elimizde vektör yok; yalnızca birbirine bağlı üç
sayı var. Bu değerler ayrı ayrı \(40\), \(2\), \(15\) biçiminde
tutulabilir; fakat o zaman hangi sayının neyi gösterdiği ve bunların bir
arada tek bir nesne oluşturduğu ayrıca belirtilmelidir.

Koordinat vektörü bu üç değeri sabit bir sırayla tek bir matematiksel
nesnede toplar. Burada önemli olan yalnız üç sayının alt alta yazılması
değildir; bu sayıların sabit bir sırayla ve tek nesne olarak ele
alınmasıdır. Böylece tüm gruba tek işlemle davranabiliriz: bütün
miktarları iki katına çıkarmak istersek tek tek üç işlem yerine \(2v\)
yazmak yeterlidir.

\begin{center}\rule{0.5\linewidth}{0.5pt}\end{center}

\subsection{Koordinat Vektörü ve
Bileşenler}\label{koordinat-vektuxf6ruxfc-ve-bileux15fenler}

\[
v=\begin{bmatrix}v_1\\v_2\\\vdots\\v_n\end{bmatrix}\in\mathbb{R}^n
\]

\begin{itemize}
\tightlist
\item
  \(v\): bütün vektör
\item
  \(v_i\): \(i\). bileşen, bir \textbf{skaler}
\item
  \(n\): bileşen sayısı
\item
  İndis konumu belirtir
\end{itemize}

\(n\) bileşenli bir koordinat vektörü bileşenleriyle gösterilir. \(v_i\)
ifadesi ``\(v\) vektörünün \(i\). bileşeni'' demektir; buradaki indis
yalnız bir etiket değil, konum bilgisidir. Örneğin ürün örneğinde
\(v_1\) fiyatı, \(v_2\) ağırlığı, \(v_3\) stok miktarını temsil eder.

Bir koordinat vektörünün bileşenleri skalerdir: \(v\in\mathbb{R}^n\) ise
her \(v_i\) gerçek bir sayıdır. Bu derste ``vektör'' dediğimizde
şimdilik \(\mathbb{R}^n\) içindeki koordinat vektörlerini kastediyoruz.
Bu sınırlama bilinçli bir başlangıçtır; vektörün daha genel tanımı
ileride vektör uzaylarında gelecek ve orada polinomların,
fonksiyonların, hatta matrislerin de vektör rolü üstlenebildiğini
göreceğiz.

\begin{center}\rule{0.5\linewidth}{0.5pt}\end{center}

\subsection{Bileşenlerin Sırası Yapının
Parçasıdır}\label{bileux15fenlerin-sux131rasux131-yapux131nux131n-paruxe7asux131dux131r}

\[
\begin{bmatrix}2\\-1\\4\end{bmatrix}
\neq
\begin{bmatrix}-1\\2\\4\end{bmatrix}
\]

Sıcaklık--basınç--debi okuması:

\[
\begin{bmatrix}80\\25\\120\end{bmatrix}
\neq
\begin{bmatrix}25\\80\\120\end{bmatrix}
\]

Bir vektörde yalnızca hangi değerlerin bulunduğu değil, bu değerlerin
hangi sırada bulunduğu da önemlidir. Yukarıdaki iki vektör aynı üç
sayıyı içerir; fakat birinci ve ikinci konumlar değiştiği için aynı
vektör değildir.

Bileşenler belirli özellikleri temsil ettiğinde bu fark daha açık
görülür. Bileşenleri sırasıyla sıcaklık, basınç ve debi olan bir
vektörde ilk iki değer yer değiştirirse, sayılar aynı kalsa bile artık
aynı fiziksel bilgi temsil edilmez: \(80\) derece ve \(25\) bar yerine
\(25\) derece ve \(80\) bar okunur. Sensör sırası, ürün sırası veya
bilinmeyen sırası değişirse problem de değişir.

\begin{center}\rule{0.5\linewidth}{0.5pt}\end{center}

\subsection{Vektörlerle Temsil Edilen
Nicelikler}\label{vektuxf6rlerle-temsil-edilen-nicelikler}

\begin{longtable}[]{@{}ll@{}}
\toprule\noalign{}
Bağlam & Vektör \\
\midrule\noalign{}
\endhead
\bottomrule\noalign{}
\endlastfoot
Konum & \((x,y,z)^T\) \\
Hız & \((v_x,v_y,v_z)^T\) \\
Sistem ölçümü & (sıcaklık, basınç, debi)\(^T\) \\
Sınav puanları & \((p_1,\ldots,p_k)^T\) \\
Veri örneği & özellik vektörü \((x_1,\ldots,x_n)^T\) \\
\end{longtable}

Vektörün bileşenlerinin mutlaka fiziksel koordinatlar olması gerekmez.
Bir cismin uzaydaki konumu, üç yöndeki hız bileşenleri, bir sistemin
sıcaklık-basınç-debi ölçümleri, bir öğrencinin farklı sınavlardan aldığı
puanlar veya bir veri örneğinin farklı özellikleri aynı biçimde vektörle
temsil edilebilir.

Bir görüntünün renk değerleri, bir ekonomik sistemin göstergeleri, bir
sensör grubundan alınan ölçümler ya da bir optimizasyon probleminin
değişkenleri de uygun bir sırayla vektöre yerleştirilebilir. Bu
örneklerin ortak noktası, birbiriyle ilişkili birden fazla skaler
değerin anlamlı ve sabit bir sıra içinde tek bir nesne olarak ele
alınmasıdır. Vektörün bileşenleri aynı fiziksel birimde olmak zorunda da
değildir; fiyat, ağırlık ve stok birlikte taşınabilir. Ancak böyle bir
durumda yapılacak işlemin problem bağlamında anlamlı olup olmadığı
ayrıca kontrol edilmelidir.

\begin{center}\rule{0.5\linewidth}{0.5pt}\end{center}

\subsection{Vektörün Ait Olduğu
Uzay}\label{vektuxf6ruxfcn-ait-olduux11fu-uzay}

\[
v\in\mathbb{R}^n
\quad\Longleftrightarrow\quad
v\ \text{'nin } n \text{ gerçek bileşeni var}
\]

\[
u=\begin{bmatrix}3\\-2\end{bmatrix}\in\mathbb{R}^2,
\qquad
w=\begin{bmatrix}1\\0\\4\\-3\end{bmatrix}\in\mathbb{R}^4
\]

Bir vektörün kaç bileşenden oluştuğunu ve bileşenlerinin hangi skaler
cisminden geldiğini belirtmek için ait olduğu uzay yazılır.
\(v\in\mathbb{R}^3\) ifadesi, \(v\)'nin üç gerçek sayı bileşeninden
oluştuğunu söyler.

Günlük matematik dilinde \(w\) için ``dört boyutlu vektör'' ifadesi de
yaygındır. Daha dikkatli bir ifadeyle \(w\), dört bileşenli bir vektör
ve dört boyutlu \(\mathbb{R}^4\) uzayının bir elemanıdır. Koordinat
uzaylarında bileşen sayısı ile uzayın boyutu aynı sayı olduğu için bu
iki kullanım çoğu zaman karışmaz; daha genel vektör uzaylarında ise
boyut, uzayın ayrıca tanımlanan bir özelliği olarak ele alınacaktır.

\begin{center}\rule{0.5\linewidth}{0.5pt}\end{center}

\subsection{Satır mı, Sütun mu?}\label{satux131r-mux131-suxfctun-mu}

\[
v=\begin{bmatrix}2\\-1\\4\end{bmatrix}\ (3\times1),
\qquad
v^T=\begin{bmatrix}2&-1&4\end{bmatrix}\ (1\times3)
\]

Bu derste aksi belirtilmedikçe vektörler \textbf{sütun vektörü}dür.

Aynı bileşenler iki yazımda da bulunur; fakat şekil değişir. Sütun
vektörü \(3\times1\), satır vektörü \(1\times3\) boyutundadır ve matris
işlemlerinde bu iki şekil aynı davranmaz.

Sütun tercihi keyfî değildir. İleride \(Ax\) yazdığımızda \(A\) bir
matris, \(x\) bir sütun vektörü olacak; \(A\in\mathbb{R}^{m\times n}\)
ve \(x\in\mathbb{R}^n\) ise \(Ax\) çarpımı \(m\) bileşenli bir sütun
vektörü verir. Vektörleri sütun olarak yazmak, \(Ax=b\) denklem sistemi
notasyonunu tutarlı kılar. Satır ve sütun vektörleri arasındaki ilişkiyi
transpoz işlemiyle daha sonra açıkça kuracağız.

\begin{center}\rule{0.5\linewidth}{0.5pt}\end{center}

\subsection{Vektör Eşitliği}\label{vektuxf6r-eux15fitliux11fi}

İki koordinat vektörü ancak

\begin{enumerate}
\def\labelenumi{\arabic{enumi}.}
\tightlist
\item
  aynı sayıda bileşene sahipse,
\item
  karşılık gelen bütün bileşenleri eşitse
\end{enumerate}

eşittir.

\[
\begin{bmatrix}2a-1\\b+3\end{bmatrix}
=
\begin{bmatrix}5\\7\end{bmatrix}
\Longrightarrow
a=3,\quad b=4
\]

Vektör eşitliği, tek bir yazımı bileşen düzeyindeki skaler eşitliklere
açar. Örnekteki eşitlik aslında \(2a-1=5\) ve \(b+3=7\) denklemlerini
aynı anda söylemektir; buradan \(a=3\) ve \(b=4\) bulunur.

Bu mekanizma, ileride \(Ax=b\) ifadesinin birden fazla lineer denklemi
tek bir matris yazımı içinde nasıl taşıdığını anlamaya hazırlık sağlar.
Aynı bileşenleri farklı sırada taşıyan iki vektör eşit değildir; konum
bilgisi eşitliğin parçasıdır ve bileşenlerin küme olarak aynı olması
yeterli olmaz.

\begin{center}\rule{0.5\linewidth}{0.5pt}\end{center}

\subsection{Skaler Neden Ayrı Bir
Ad?}\label{skaler-neden-ayrux131-bir-ad}

Vektörleri ölçeklemek için kullanılan sayılara \textbf{skaler} denir.

\[
3
\begin{bmatrix}
2\\
-1
\end{bmatrix}
\]

Buradaki \(3\), ilk bakışta yalnızca bir \textbf{katsayı} gibi görünür.

\begin{quote}
Peki neden ayrıca \textbf{skaler} diyoruz?
\end{quote}

Skaler kavramıyla ilk kez karşılaştığımızda ayrı bir ad kullanmak
gereksiz görünebilir. Çünkü şu ana kadar gördüğümüz örneklerde skalerler
\(2\), \(-3\) veya \(\frac12\) gibi bildiğimiz sayılardır. Bir vektörün
önündeki \(3\) için yalnızca ``katsayı'' demek de hesabı yapmak için
yeterlidir.

Bu aşamada gerçekten de skaleri basit bir katsayı gibi düşünebiliriz.
Ancak ileride vektör uzaylarını ele aldığımızda, bir vektörü hangi
elemanlarla çarpabileceğimiz yapının tanımının bir parçası olacaktır.
Yani her durumda kullanılabilecek katsayılar gelişigüzel seçilmez.

Bu nedenle ``skaler'' adını baştan kullanıyoruz. Şimdilik basit bir
katsayı gibi görünse de kavramın arkasında ileride önemli hale gelecek
teknik bir ayrım vardır. Bunun ilk işaretini, skalerlerin hangi sayı
sisteminden seçildiğine bakarak görebiliriz.

\begin{center}\rule{0.5\linewidth}{0.5pt}\end{center}

\subsection{Skaler Cismi}\label{skaler-cismi}

\textbf{Meraklısına not:} Skalerler belirli bir \textbf{cisimden}
seçilir:

\[
\mathbb{R}
\qquad\text{veya}\qquad
\mathbb{C}
\]

\[
2-3i:
\quad
\mathbb{R}\text{ üzerinde skaler değil},
\qquad
\mathbb{C}\text{ üzerinde geçerli skaler}
\]

Peki:

\begin{itemize}
\tightlist
\item
  \textbf{Cisim nedir?}
\item
  \textbf{Bunu bilmek neden gerekli?}
\end{itemize}

Bir önceki slaytta, skalerin şimdilik sıradan bir katsayı gibi
göründüğünü söyledik. Buradaki teknik ayrım, hangi elemanların bu
katsayı rolünü üstlenebileceğiyle ilgilidir.

Skalerler belirli bir cisimden seçilir. Bir cisim, kabaca toplama,
çıkarma, çarpma ve sıfıra bölme dışında bölme işlemlerinin alıştığımız
kurallarla yapılabildiği bir sayı sistemidir. Gerçek sayılar
\(\mathbb{R}\) ve karmaşık sayılar \(\mathbb{C}\) en sık
karşılaşacağımız örneklerdir. Ayrıntılı tanımına burada ihtiyacımız
olmayacak.

Örneğin gerçek sayılar üzerinde çalışıyorsak \(2-3i\) bir skaler olarak
kullanılamaz. Karmaşık sayılar üzerinde çalışıyorsak kullanılabilir.
Başlangıçta bu ayrım çok önemli görünmeyebilir; çünkü çoğunlukla gerçek
sayılarla çalışacağız. Ancak vektör uzaylarına geçtiğimizde hangi cismin
kullanıldığı, o uzayın tanımının doğrudan bir parçası olacaktır.

Dolayısıyla burada amaç cisim kavramını ayrıntılı olarak öğrenmek değil,
``skaler'' sözcüğünün neden yalnızca ``katsayı'' yerine kullanıldığını
görmek: bugün basit görünen bu ayrım, ileride matematiksel yapının
kendisini belirleyecektir.

\begin{center}\rule{0.5\linewidth}{0.5pt}\end{center}

\subsection{Sıfır, Bir ve Negatif
Skalerler}\label{sux131fux131r-bir-ve-negatif-skalerler}

\[
1v=v
\qquad\text{(etkisiz skaler)}
\]

\[
0v=0
\qquad\text{(sıfır vektörüne götürür)}
\]

\[
(-1)v=-v
\qquad\text{(yönü ters çevirir)}
\]

Bazı skalerlerin ölçekleme açısından özel rolleri vardır. \(1\) etkisiz
skalerdir: bir vektörü \(1\) ile çarpmak onu değiştirmez. Bu,
sayılardaki \(1\cdot a=a\) kuralının vektör karşılığıdır.

\(0\) skaleri her vektörü ilgili uzayın sıfır vektörüne götürür; dikkat
edilmesi gereken nokta, soldaki \(0\)'ın bir skaler, sağdaki \(0\)'ın
ise bir vektör olmasıdır. Negatif bir skaler, gerçek vektörlerin
geometrik yorumunda ölçeklemenin yanında yönün tersine dönmesine neden
olur; özel olarak \((-1)v\) çarpımı \(v\)'nin toplamsal tersi olan
\(-v\) vektörünü verir. Bu üç sonuç, birazdan göreceğimiz cebirsel
özelliklerden de türetilebilir.

\begin{center}\rule{0.5\linewidth}{0.5pt}\end{center}

\subsection{Hangileri Skalerdir?}\label{hangileri-skalerdir}

\(\mathbb{R}\) üzerinde çalışan bir uzayda:

\[
-3,\qquad
\sqrt{2},\qquad
2+i,\qquad
\begin{bmatrix}4\\1\end{bmatrix},\qquad
\begin{bmatrix}5\end{bmatrix}
\]

İlk ikisi, \(-3\) ve \(\sqrt2\), gerçek sayı oldukları için bu uzayda
skalerdir. Üçüncüsü, \(2+i\), bir karmaşık sayıdır; \(\mathbb{R}\)
üzerinde çalışan bu uzayın skaleri değildir. Aynı sayı, \(\mathbb{C}\)
üzerinde tanımlı bir uzayda geçerli bir skaler olurdu --- yani ``skaler
olmak'' nesnenin kendisine değil, çalışılan uzaya bağlıdır.

Dördüncü nesne iki bileşenli bir vektördür, skaler değildir. Beşinci
nesne tek bileşenli olsa da vektör olarak yazılmıştır; bu nedenle skaler
\(5\) ile aynı matematiksel nesne değildir. Bu son ayrım özellikle
önemlidir ve bir sonraki slaytta ayrıntılandırılacaktır.

\begin{center}\rule{0.5\linewidth}{0.5pt}\end{center}

\subsection{Aynı Sayı, Farklı
Nesne}\label{aynux131-sayux131-farklux131-nesne}

\begin{longtable}[]{@{}ll@{}}
\toprule\noalign{}
Yazım & Bu dersteki okuma \\
\midrule\noalign{}
\endhead
\bottomrule\noalign{}
\endlastfoot
\(5\) & skaler \\
\(\begin{bmatrix}5\end{bmatrix}\) & tek bileşenli sütun vektörü \\
\([5]\) & bağlama göre \(1\times1\) matris \\
\end{longtable}

\begin{quote}
İçerilen sayı aynı olabilir; matematiksel rol aynı olmak zorunda
değildir.
\end{quote}

\(5\) ile \(\begin{bmatrix}5\end{bmatrix}\) aynı değeri içeriyor gibi
görünür; fakat işlemde aynı tür nesne olarak davranmazlar. \(5\) bir
katsayı olabilir, \(\begin{bmatrix}5\end{bmatrix}\) ise tek bileşenli
bir sütun vektörüdür ve bir şekli vardır.

Bu ayrım ileride boyut kontrollerinde önem kazanır. \(5v\) ifadesinde
\(5\) vektörü ölçekler; buna karşılık \(\begin{bmatrix}5\end{bmatrix}v\)
aynı şey değildir, çünkü soldaki nesne bir sayı değil şekli olan başka
bir nesnedir. Dolayısıyla ``içinde tek sayı bulunan her şey skalerdir''
düşüncesi matris çarpımında ve denklem sistemlerinde tür hatalarına yol
açar. Bir işlemi doğru yorumlamak için önce nesnenin türüne bakılır:
skaler mi, vektör mü, matris mi?

\begin{center}\rule{0.5\linewidth}{0.5pt}\end{center}

\subsection{Skalerin İki Rolü}\label{skalerin-iki-roluxfc}

\textbf{1. Bileşen olarak:}

\[
v=\begin{bmatrix}2\\-3\\5\end{bmatrix},
\qquad
A=\begin{bmatrix}1&4\\-2&3\end{bmatrix}
\]

\textbf{2. Katsayı olarak:}

\[
\alpha v,\qquad \alpha A
\]

Skalerler lineer cebirde iki temel biçimde karşımıza çıkar. Birincisi
bileşen rolüdür: bir vektörün bileşenleri ve bir matrisin girdileri
skalerdir. Örnekteki \(2\), \(-3\), \(5\) değerleri ile \(1\), \(4\),
\(-2\), \(3\) değerlerinin her biri birer skalerdir.

İkincisi katsayı rolüdür: bir skaler, bir vektör veya matrisin önünde
onu ölçekleyen çarpan olarak kullanılır. \(3v\) ifadesindeki \(3\) ve
\(-2A\) ifadesindeki \(-2\) bu roldedir. Aynı sayı iki rolde de
görünebilir; önemli olan, hangi rolde kullanıldığını takip etmektir.
Vektör ve matrisler skalerlerden oluşsalar bile kendi başlarına farklı
türde matematiksel nesnelerdir.

\begin{center}\rule{0.5\linewidth}{0.5pt}\end{center}

\subsection{Cebirsel Nesne, Geometrik
Temsil}\label{cebirsel-nesne-geometrik-temsil}

\[
v=\begin{bmatrix}3\\2\end{bmatrix}
\]

\begin{itemize}
\tightlist
\item
  Cebirsel: iki bileşenli koordinat vektörü
\item
  Geometrik: yatayda \(3\), düşeyde \(2\) birim yer değiştirme
\end{itemize}

\begin{quote}
Ok, vektörün bir \textbf{temsilidir} --- genel tanımı değil.
\end{quote}

Bir koordinat vektörü cebirsel olarak bileşenleriyle ifade edilir. Aynı
vektör geometrik olarak yatayda \(3\), düşeyde \(2\) birimlik yer
değiştirmeyi gösteren yönlü bir doğru parçasıyla, yani bir okla temsil
edilebilir. Bu iki okuma birbirini tamamlar: cebirsel gösterim vektörü
bileşenleriyle ifade eder, geometrik gösterim uzaydaki etkisini
görselleştirir.

Ancak geometrik ok, vektör kavramının genel tanımı değildir; yalnızca
\(\mathbb{R}^2\) ve \(\mathbb{R}^3\) gibi görselleştirilebilen uzaylarda
kullanılan bir temsil biçimidir. Bu ayrımı erken yapmak önemlidir:
ileride fonksiyonlar veya matrisler vektör olarak ele alındığında onlara
ok çizmek mümkün olmayabilir; bu, onların vektör olmadığı anlamına
gelmez.

\begin{center}\rule{0.5\linewidth}{0.5pt}\end{center}

\subsection{Nokta ve Vektör Aynı Şey
Değildir}\label{nokta-ve-vektuxf6r-aynux131-ux15fey-deux11fildir}

\[
P=(3,2)
\quad\text{nokta: bir \textbf{konum}}
\]

\[
v=\begin{bmatrix}3\\2\end{bmatrix}
\quad\text{vektör: bir \textbf{yer değiştirme}}
\]

\[
\overrightarrow{OP}=\begin{bmatrix}3\\2\end{bmatrix}
\quad (O=(0,0)\ \text{konum vektörü})
\]

Düzlemde \(P=(3,2)\) ifadesi bir noktayı gösterir ve bu noktanın belirli
bir konumu vardır. Buna karşılık aynı sayıları taşıyan \(v\) bir
vektördür ve geometrik bağlamda bir yönlü yer değiştirmeyi temsil eder.
İki gösterimde aynı sayılar kullanılsa da matematiksel rolleri
farklıdır.

Orijin \(O=(0,0)\) olmak üzere \(P\) noktasına karşılık gelen
\(\overrightarrow{OP}\) vektörüne \(P\) noktasının konum vektörü denir.
Bu özel durumda noktanın koordinatları ile konum vektörünün bileşenleri
aynı sayılardır; bu, iki kavramın karıştırılmasının başlıca nedenidir.
Ancak sayısal örtüşme rol özdeşliği anlamına gelmez: nokta nerede
olunduğunu, vektör ne kadar ve hangi yönde gidildiğini söyler.

\begin{center}\rule{0.5\linewidth}{0.5pt}\end{center}

\subsection{Serbest Vektör}\label{serbest-vektuxf6r}

\[
A=(1,1),\ B=(4,3)
\ \Rightarrow\
\overrightarrow{AB}=\begin{bmatrix}3\\2\end{bmatrix}
\]

\[
C=(-2,4),\ D=(1,6)
\ \Rightarrow\
\overrightarrow{CD}=\begin{bmatrix}3\\2\end{bmatrix}
\]

\[
\boxed{\overrightarrow{AB}=\overrightarrow{CD}}
\]

Bir vektör geometrik olarak bir okla gösterildiğinde, okun başlangıç
noktası vektörün kendisinin bir parçası değildir. \(A=(1,1)\) ile
\(B=(4,3)\) arasındaki yer değiştirme \((3,2)^T\)'dir. Tamamen farklı
bir bölgedeki \(C=(-2,4)\) ile \(D=(1,6)\) noktaları arasındaki yer
değiştirme de aynı vektördür.

İki ok düzlemin farklı yerlerinde çizilmiş olsa da aynı yönlü yer
değiştirmeyi temsil eder; yönleri ve uzunlukları aynıdır. Başlangıç
noktasından bağımsız düşünülen bu vektörlere serbest vektör denir. Bu
nedenle bir vektörü orijinden başlatarak çizmek yalnızca standart ve
kullanışlı bir gösterim tercihidir. Serbest vektör yorumu, ileride
vektör toplamada ikinci vektörü birincinin ucuna paralel
taşıyabilmemizin de dayanağıdır.

\begin{center}\rule{0.5\linewidth}{0.5pt}\end{center}

\subsection{Vektör--Skaler Ayrımı Daha
Geneldir}\label{vektuxf6rskaler-ayrux131mux131-daha-geneldir}

Şimdilik çoğunlukla

\[
v=
\begin{bmatrix}
v_1\\
\vdots\\
v_n
\end{bmatrix},
\qquad
\alpha\in\mathbb{R}
\]

biçiminde çalışacağız.

Ancak genel olarak:

\begin{quote}
\textbf{Vektör}, üzerinde çalışılan vektör uzayının bir elemanıdır.\\
\textbf{Skaler}, bu uzayın skaler cisminden gelir.
\end{quote}

Bu nedenle uygun bir vektör uzayında

\[
p(x)=x^2+1
\]

de bir vektör olabilir ve

\[
3p(x)
\]

ifadesindeki \(3\) yine skalerdir.

Başlangıçta vektörleri birden fazla sayıyı bir arada tutan nesneler,
skalerleri ise tek sayılar olarak gördük. Şu an çalışacağımız koordinat
vektörleri için bu oldukça kullanışlı bir ayrımdır. Ancak bunu vektör ve
skalerin genel tanımı olarak düşünmemek gerekir.

İleride vektör uzaylarına geçtiğimizde vektör kavramının sayı
sütunlarından daha genel olduğunu göreceğiz. Örneğin uygun bir vektör
uzayında bir polinom veya matris de vektör olabilir. Bu durumda onu
çarpan elemanlar yine skalerlerdir.

Burada bu yapıyı ayrıntılı olarak incelemiyoruz. Bu noktayı yalnızca
başlangıçta kurduğumuz ``vektör çok sayıdır, skaler tek sayıdır''
sezgisinin çalışma kolaylığı sağlayan geçici bir ayrım olduğunu
belirtmek için söylüyoruz. Genel ayrımı vektör uzayları konusunda
kuracağız.

\begin{center}\rule{0.5\linewidth}{0.5pt}\end{center}

\subsection{Sık Yapılan Hatalar}\label{sux131k-yapux131lan-hatalar}

\begin{enumerate}
\def\labelenumi{\arabic{enumi}.}
\tightlist
\item
  ``İçinde tek sayı bulunan her şey skalerdir'' sanmak:
  \(\begin{bmatrix}5\end{bmatrix}\neq5\).
\item
  Skaler olmayı nesnenin kendisine bağlamak --- çalışılan uzaya bağlıdır
  (\(2+i\)).
\item
  Bileşenlerin sırasını önemsiz saymak.
\item
  Noktayı vektörle özdeşleştirmek (konum ≠ yer değiştirme).
\item
  Geometrik oku vektörün \textbf{tanımı} sanmak.
\end{enumerate}

Birinci hata nesne türü ile içerdiği sayıyı karıştırır; tek bileşenli
bir sütun vektörünün şekli vardır ve skalerle aynı nesne değildir. Bu
hata matris çarpımında ve denklem sistemlerinde boyut hatalarına yol
açar. İkinci hata skaler olmayı mutlak bir özellik sanmaktır; \(2+i\)
gerçek bir uzayda skaler değilken karmaşık bir uzayda geçerli skalerdir.

Üçüncü hata bileşen sırasını gözden kaçırmaktır; sıra yapının parçasıdır
ve değiştiğinde vektör de değişir. Dördüncü hata nokta ile vektörü
özdeşleştirmektir; konum vektörü durumunda sayılar örtüşse de roller
farklıdır. Beşinci hata oku tanım sanmaktır; ok yalnız \(\mathbb{R}^2\)
ve \(\mathbb{R}^3\) için bir temsildir ve polinom ya da matris
vektörlerine çizilemez --- bu onların vektör olmadığı anlamına gelmez.

\begin{center}\rule{0.5\linewidth}{0.5pt}\end{center}

\subsection{Sonraki Adım: İki Temel
İşlem}\label{sonraki-adux131m-iki-temel-iux15flem}

Nesneleri ayırdık: \textbf{skaler} ve \textbf{koordinat vektörü}.

\[
u+v
\qquad\text{ve}\qquad
\alpha v
\]

\begin{quote}
Sırada: bu iki işlemi kurmak ve \textbf{lineer birleşim} içinde
birleştirmek.
\end{quote}

Bu derste iki nesne türünü ayırdık: vektör uzayının elemanı olan
koordinat vektörleri ve bu uzayın skaler cisminden gelen, vektörleri
ölçekleyen skalerler. Bileşenlerin sırasının yapının parçası olduğunu,
eşitliğin bileşen düzeyinde tanımlandığını ve geometrik okun yalnız bir
temsil olduğunu gördük.

Bir sonraki adımda bu nesneler üzerinde iki temel işlem kurulacak:
vektör toplama ve skalerle çarpma. Çıkarma ayrı bir temel işlem değil,
bu ikisinin birlikte kullanılmasıyla elde edilecek. Sonunda aynı iki
işlem lineer birleşim kalıbında birleşecek; bu kalıp, dersin geri
kalanının ana dili olacaktır.




\end{document}
